% Section 11 — Conclusion (charter §4.11; issue T4-1)
\section{Conclusion}
\label{sec:conclusion}

We set out to explain, rather than merely report, the central null of
\texttt{preprint-v1}, and to do so without converting a null into a manufactured
positive. The explanation is two-halved. The negative half is a theorem: under
information parity the data-processing inequality closes at the optimum, and the
Bayes risk is format-invariant (\cref{thm:irrelevance}), so the Lumen T1 null is
the \emph{predicted} information-limit outcome and not an anomaly to be explained
away (guardrail~G1). When parity fails, the failure is not catastrophic but
graded: the inflation of Bayes risk is the lost information exactly under log
loss and within a $\sqrt{\cdot}$ envelope otherwise (\cref{cor:ratedist-gap}).
The positive half localizes, and bounds, the only place a real effect can live.
A bounded learner is not the Bayes posterior, and the room between them is a
capacity budget with no information content (\cref{thm:bounded-gap},
\cref{cor:size-gap}); that budget is non-empty only in the finite-resource,
capacity-separated window the theory names (\cref{prop:regime}). The two go/no-go
refinements sharpen the picture without changing it: operational parity is a
retraction whose defect hides from both the mutual information and the Bayes risk
(\cref{thm:retraction}), and that same defect, once metrized, dominates the swap
gap of any fixed bounded predictor (\cref{thm:defect-risk-gap}).

Read together, the three frozen Lumen cells occupy three distinct regimes ---
T1 at the closed information channel, T3 at sub-frontier saturation, T2 in the
permitted-but-unidentified activation window (\cref{sec:anchoring}) --- and we
claim of this mapping only consistency, never confirmation
(\cref{rem:anchor-limits}). The theory earns its keep not by retrodiction but by
risk: the five predictions of \cref{sec:predictions} attach a sign-aware trend,
a prohibition, or a bound to a concrete knob of the apparatus, each with an
explicit falsifier (\cref{tab:pred-gates}). A parity-verified effect that
\emph{grows} as the model recedes from its frontier would falsify the thesis; a
clean null in the frontier regime would corroborate the negative half; a
confirmed monotone effect would scope the product to the bounded regime the
theory marks out. None of these outcomes leaves the theory ``always consistent,''
which is the point.

The honest residue is small but real, and \cref{sec:limitations} collects it.
The irrelevance theorem lives at two idealizations a frozen LLM does not meet;
the bounded-learner objects map onto a single model by analogy; the T1/T3
distinction rests on an empirical surmise about task difficulty that is
conjectural, not proved (guardrail~G6). What survives all of this is a single,
falsifiable claim: format cannot buy information that is not there, so any
genuine code-representation effect on a bounded model is a capacity phenomenon of
bounded size, appearing only near the model's frontier --- and that is a
statement the apparatus can be pointed at and made to test.
